\section{Wet artificial life: construction, replication, and integration}
\label{sec:wet-life}

Wet artificial life investigates life-like organization in biochemical systems.
Bottom-up work constructs selected functions from molecular components;
top-down work simplifies or redesigns existing cells. Their achievements concern
distinct functions: compartment growth, molecular replication, cellular
reproduction, or adaptation. Membrane division does not demonstrate genome
copying; genome synthesis does not establish construction of every cellular
component from nonliving materials
\citep{wetZhu2009,wetGibson2010,wetHutchison2016}.

\subsection{Compartments: reproduction without a complete genetic system}

Zhu and Szostak demonstrated coupled growth and division of model protocell
membranes. Fatty-acid vesicles supplied with membrane material developed
elongated shapes that divided under modest fluid shear. Repeated cycles were
possible, and encapsulated RNA was redistributed into daughter compartments
without the substantial leakage associated with their comparison method
\citep{wetZhu2009}. The result establishes a physical route linking boundary
growth, division, and retention of contents, rather than treating division as an
unrelated operation imposed after growth.

The distinction between distributing RNA and replicating it is essential.
The study did not demonstrate an internally replicating genome coupled to the
membrane cycle. Its Discussion explicitly identifies that integration as a
further requirement for a more complete protocell. Membrane precursors and
mechanical agitation remained externally supplied. The experiment isolates a
tractable physical problem; its success does not establish cellular autonomy
or sustained adaptive evolution \citep{wetZhu2009}.

\subsection{Chemical replication and limited evolutionary dynamics}

Lincoln and Joyce constructed RNA enzymes that catalysed one another's
production from supplied component substrates. The cross-catalytic system
supported exponential amplification without proteins in the replication
reaction. When related replicators competed for shared substrates, recombinant
variants arose and became abundant \citep{wetLincoln2009}. This provides a
chemical example in which replication machinery participates in the system
being reproduced, and in which heritable differences affect competitive
outcomes. It goes beyond observing a persistent structure or merely counting
more molecules.

Here, ``self-sustained'' describes a catalytic cycle conditional on substrate
replenishment, rather than synthesis of those substrates, membrane maintenance,
or environmental control. The enzymes had also undergone prior experimental
optimization. The authors characterize their genetic system as restricted and
identify the emergence of new encoded functions beyond replication as an open
challenge. Thus the work demonstrates a bounded form of chemical evolution,
not open-ended functional innovation or a complete account of life's origin
\citep{wetLincoln2009}. It offers a useful comparison for digital replicators:
the source of building blocks and the allowed variation determine which
evolutionary claims the observed competition can support.

\subsection{Synthetic genomes and top-down minimal cells}

Gibson and colleagues demonstrated cells controlled by a chemically synthesized
bacterial chromosome after its introduction into a recipient cell. Descendants
reproduced with the synthetic genome and exhibited the expected phenotypic
properties \citep{wetGibson2010}. The achievement concerns genome construction
and control of an existing cellular system. The paper explicitly states that
the recipient cytoplasm was not synthetic. Inherited cellular machinery matters
to initial establishment, although its original molecules are progressively
replaced in descendants.
Genome-directed renewal and construction of an entire cell from nonliving
components remain distinct experimental claims.

Hutchison and colleagues subsequently minimized a synthetic bacterial genome
through iterative design and testing. Their initial design was not viable;
retaining genes needed for robust growth helped produce a viable reduced
genome, which still included genes of unknown biological function
\citep{wetHutchison2016}.\footnote{The account of this study is limited to
the authors' published abstract and institutional publication record;
the detailed full text was not inspected in this review pass.}
This is an empirical approach to minimality: successful reproduction tests
whether a reduced system remains viable under its experimental conditions.
The unresolved gene functions also distinguish a small working cell from a
fully understood cell. Genome size alone supplies neither a complete causal
explanation nor a universal lower bound for every environment or lineage.

\subsection{Minimality does not remove the capacity to adapt}

A later experiment directly examined evolution in a genomically minimized
cell. Moger-Reischer and colleagues compared minimal and non-minimal strains
using mutation measurements, laboratory evolution, fitness assays, and genome
analysis. The minimized strain recovered fitness lost through streamlining,
while some phenotypic responses depended on genomic context
\citep{wetMogerReischer2023}. This demonstrates that severe reduction need not
eliminate adaptive potential. It also changes the question from whether a
designed genome functions initially to how its descendants respond to selection.

The cells inherited biological organization and received externally maintained
culture conditions and supplies. Adaptation over the observation period does
not establish spontaneous origin, environmental independence, or indefinitely
expanding complexity \citep{wetMogerReischer2023}.

\subsection{The integration problem and its relevance to digital models}

Together, these studies identify a methodological tradeoff. Isolated bottom-up
systems make particular mechanisms easier to attribute, while retaining fewer
integrated cellular functions. Top-down cells reproduce and adapt as complete
organisms, but inherit machinery and organizational history that the experiment
does not reconstruct. The comparison is therefore about which dependencies
have been explained, rather than which demonstration looks most life-like.
An integration claim must connect boundary maintenance, heritable replication,
resource acquisition, and reproduction within the same continuing system.

For soft artificial life, the transferable lesson is experimental accounting,
not a requirement to imitate cellular chemistry. A digital study should state
which processes its programs perform, which resources the environment supplies,
and which lifecycle transitions the simulator grants. Distinguishing a working
replication mechanism from establishment, adaptation, and continuing innovation
helps prevent a successful component demonstration from carrying a broader
claim than its evidence supports.
